\begin{register}{H}{MEM\_MONITOR\_LOG\_SETTING\_REG}{0x0000}\label{regdesc:MEMMONITORLOGSETTINGREG}
\regfield{MEM\_MONITOR\_LOG\_DMA\_1\_ENA}{8}{24}{{0}}%
\regfield{MEM\_MONITOR\_LOG\_DMA\_0\_ENA}{8}{16}{{0}}%
\regfield{MEM\_MONITOR\_LOG\_CORE\_ENA}{8}{8}{{0}}%
\regfield{(reserved)}{3}{5}{000}%
\regfield{MEM\_MONITOR\_LOG\_MEM\_LOOP\_ENABLE}{1}{4}{{1}}%
\regfield{MEM\_MONITOR\_LOG\_MODE}{4}{0}{{0}}%
\reglabel{Reset}\regnewline%
\begin{regdesc}\begin{reglist}
\label{fielddesc:MEMMONITORLOGMODE}\item [MEM\_MONITOR\_LOG\_MODE] 配置监测模式。\\
1：使能写监测\\
2：使能 word 监测\\
4：使能 halfword 监测\\
8：使能 byte 监测\\
其他值：无效\\
(R/W)
\label{fielddesc:MEMMONITORLOGMEMLOOPENABLE}\item [MEM\_MONITOR\_LOG\_MEM\_LOOP\_ENABLE] 配置写内存的存储模式。\\
0：非 loop 模式\\
1：loop 模式\\
(R/W)
\label{fielddesc:MEMMONITORLOGCOREENA}\item [MEM\_MONITOR\_LOG\_CORE\_ENA] 配置是否使能 HP CPU 总线访问记录。\\
0：不使能\\
1：使能\\
其他值：无效\\
(R/W)

\item [见下页...]
\end{reglist}\end{regdesc}
\end{register}
\addtocounter{Regfloat}{-1}
\begin{register}{H}{MEM\_MONITOR\_LOG\_SETTING\_REG}{0x0000}
\begin{regdesc}\begin{reglist}
\item [接上页...]

\label{fielddesc:MEMMONITORLOGDMA0ENA}\item [MEM\_MONITOR\_LOG\_DMA\_0\_ENA] 配置是否使能 DMA\_0 总线访问记录。\\
0：不使能\\
1：使能\\
其他值：无效\\
(R/W)
\label{fielddesc:MEMMONITORLOGDMA1ENA}\item [MEM\_MONITOR\_LOG\_DMA\_1\_ENA] 配置是否使能 DMA\_1 总线访问记录。\\
0：不使能\\
1：使能\\
其他值：无效\\
(R/W)
\end{reglist}\end{regdesc}
\end{register}


\begin{register}{H}{MEM\_MONITOR\_LOG\_SETTING1\_REG}{0x0004}\label{regdesc:MEMMONITORLOGSETTING1REG}
\regfield{(reserved)}{16}{16}{0000000000000000}%
\regfield{MEM\_MONITOR\_LOG\_DMA\_3\_ENA}{8}{8}{{0}}%
\regfield{MEM\_MONITOR\_LOG\_DMA\_2\_ENA}{8}{0}{{0}}%
\reglabel{Reset}\regnewline%
\begin{regdesc}\begin{reglist}
\label{fielddesc:MEMMONITORLOGDMA2ENA}\item [MEM\_MONITOR\_LOG\_DMA\_2\_ENA] 配置是否使能 DMA\_2 总线访问记录。\\
0：不使能\\
1：使能\\
其他值：无效\\
(R/W)
\label{fielddesc:MEMMONITORLOGDMA3ENA}\item [MEM\_MONITOR\_LOG\_DMA\_3\_ENA]配置是否使能 DMA\_3 总线访问记录。\\
0：不使能\\
1：使能\\
其他值：无效\\
(R/W)
\end{reglist}\end{regdesc}
\end{register}


\begin{register}{H}{MEM\_MONITOR\_LOG\_CHECK\_DATA\_REG}{0x0008}\label{regdesc:MEMMONITORLOGCHECKDATAREG}
\regfield{MEM\_MONITOR\_LOG\_CHECK\_DATA}{32}{0}{{0}}%
\reglabel{Reset}\regnewline%
\begin{regdesc}\begin{reglist}
\label{fielddesc:MEMMONITORLOGCHECKDATA}\item [MEM\_MONITOR\_LOG\_CHECK\_DATA] 配置总线访问监测的特殊值。(R/W)
\end{reglist}\end{regdesc}
\end{register}


\begin{register}{H}{MEM\_MONITOR\_LOG\_DATA\_MASK\_REG}{0x000C}\label{regdesc:MEMMONITORLOGDATAMASKREG}
\regfield{(reserved)}{28}{4}{0000000000000000000000000000}%
\regfield{MEM\_MONITOR\_LOG\_DATA\_MASK}{4}{0}{{0}}%
\reglabel{Reset}\regnewline%
\begin{regdesc}\begin{reglist}
\label{fielddesc:MEMMONITORLOGDATAMASK}\item[MEM\_MONITOR\_LOG\_DATA\_MASK] 配置是否屏蔽 \hyperref[regdesc:MEMMONITORLOGCHECKDATAREG]{MEM\_MONITOR\_LOG\_CHECK\_DATA\_REG} 的相应字节。可同时屏蔽多个字节。\\
bit[0]：配置是否屏蔽 \hyperref[regdesc:MEMMONITORLOGCHECKDATAREG]{MEM\_MONITOR\_LOG\_CHECK\_DATA\_REG} 最低位字节。\\
0：不屏蔽\\
1：屏蔽\\
bit[1]：配置是否屏蔽 \hyperref[regdesc:MEMMONITORLOGCHECKDATAREG]{MEM\_MONITOR\_LOG\_CHECK\_DATA\_REG} 第二个低位字节。\\
0：不屏蔽\\
1：屏蔽\\
bit[2]：配置是否屏蔽 \hyperref[regdesc:MEMMONITORLOGCHECKDATAREG]{MEM\_MONITOR\_LOG\_CHECK\_DATA\_REG} 第二个高位字节。\\
0：不屏蔽\\
1：屏蔽\\
bit[3]：配置是否屏蔽 \hyperref[regdesc:MEMMONITORLOGCHECKDATAREG]{MEM\_MONITOR\_LOG\_CHECK\_DATA\_REG} 最高位字节。\\
0：不屏蔽\\
1：屏蔽\\
(R/W)
\end{reglist}\end{regdesc}
\end{register}


\begin{register}{H}{MEM\_MONITOR\_LOG\_MIN\_REG}{0x0010}\label{regdesc:MEMMONITORLOGMINREG}
\regfield{MEM\_MONITOR\_LOG\_MIN}{32}{0}{{0}}%
\reglabel{Reset}\regnewline%
\begin{regdesc}\begin{reglist}
\label{fielddesc:MEMMONITORLOGMIN}\item [MEM\_MONITOR\_LOG\_MIN] 配置监测地址的下边界。(R/W)
\end{reglist}\end{regdesc}
\end{register}


\begin{register}{H}{MEM\_MONITOR\_LOG\_MAX\_REG}{0x0014}\label{regdesc:MEMMONITORLOGMAXREG}
\regfield{MEM\_MONITOR\_LOG\_MAX}{32}{0}{{0}}%
\reglabel{Reset}\regnewline%
\begin{regdesc}\begin{reglist}
\label{fielddesc:MEMMONITORLOGMAX}\item [MEM\_MONITOR\_LOG\_MAX] 配置监测地址的上边界。(R/W)
\end{reglist}\end{regdesc}
\end{register}


\begin{register}{H}{MEM\_MONITOR\_LOG\_MON\_ADDR\_UPDATE\_0\_REG}{0x0018}\label{regdesc:MEMMONITORLOGMONADDRUPDATE0REG}
\regfield{MEM\_MONITOR\_LOG\_MON\_ADDR\_ALL\_UPDATE}{1}{31}{{0}}%
\regfield{(reserved)}{23}{8}{00000000000000000000000}%
\regfield{MEM\_MONITOR\_LOG\_MON\_ADDR\_CORE\_UPDATE}{8}{0}{{0}}%
\reglabel{Reset}\regnewline%
\begin{regdesc}\begin{reglist}
\label{fielddesc:MEMMONITORLOGMONADDRCOREUPDATE}\item [MEM\_MONITOR\_LOG\_MON\_ADDR\_CORE\_UPDATE] 配置是否将 \hyperref[regdesc:MEMMONITORLOGMINREG]{MEM\_MONITOR\_LOG\_MIN\_REG} 到 \hyperref[regdesc:MEMMONITORLOGMAXREG]{MEM\_MONITOR\_LOG\_MAX\_REG} 的地址更新为 HP CPU 总线的监测地址空间。
\\
0：不更新\\
1：更新\\
其他值：无效\\
(WT)
\label{fielddesc:MEMMONITORLOGMONADDRALLUPDATE}\item [MEM\_MONITOR\_LOG\_MON\_ADDR\_ALL\_UPDATE] 配置是否将 \hyperref[regdesc:MEMMONITORLOGMINREG]{MEM\_MONITOR\_LOG\_MIN\_REG} 到 \hyperref[regdesc:MEMMONITORLOGMAXREG]{MEM\_MONITOR\_LOG\_MAX\_REG} 的地址更新为所有主机的监测地址空间。\\
0：不更新\\
1：更新\\
(WT)
\end{reglist}\end{regdesc}
\end{register}


\begin{register}{H}{MEM\_MONITOR\_LOG\_MON\_ADDR\_UPDATE\_1\_REG}{0x001C}\label{regdesc:MEMMONITORLOGMONADDRUPDATE1REG}
\regfield{MEM\_MONITOR\_LOG\_MON\_ADDR\_DMA\_3\_UPDATE}{8}{24}{{0}}%
\regfield{MEM\_MONITOR\_LOG\_MON\_ADDR\_DMA\_2\_UPDATE}{8}{16}{{0}}%
\regfield{MEM\_MONITOR\_LOG\_MON\_ADDR\_DMA\_1\_UPDATE}{8}{8}{{0}}%
\regfield{MEM\_MONITOR\_LOG\_MON\_ADDR\_DMA\_0\_UPDATE}{8}{0}{{0}}%
\reglabel{Reset}\regnewline%
\begin{regdesc}\begin{reglist}
\label{fielddesc:MEMMONITORLOGMONADDRDMA0UPDATE}\item [MEM\_MONITOR\_LOG\_MON\_ADDR\_DMA\_0\_UPDATE]  \hyperref[regdesc:MEMMONITORLOGMINREG]{MEM\_MONITOR\_LOG\_MIN\_REG} 到 \hyperref[regdesc:MEMMONITORLOGMAXREG]{MEM\_MONITOR\_LOG\_MAX\_REG} 的地址更新为 DMA\_0 总线的监测地址空间。\\
0：不更新\\
1：更新\\
其他值：无效\\
(WT)
\label{fielddesc:MEMMONITORLOGMONADDRDMA1UPDATE}\item [MEM\_MONITOR\_LOG\_MON\_ADDR\_DMA\_1\_UPDATE]  \hyperref[regdesc:MEMMONITORLOGMINREG]{MEM\_MONITOR\_LOG\_MIN\_REG} 到 \hyperref[regdesc:MEMMONITORLOGMAXREG]{MEM\_MONITOR\_LOG\_MAX\_REG} 的地址更新为 DMA\_1 总线的监测地址空间。\\
0：不更新\\
1：更新\\
其他值：无效\\
(WT)
\label{fielddesc:MEMMONITORLOGMONADDRDMA2UPDATE}\item [MEM\_MONITOR\_LOG\_MON\_ADDR\_DMA\_2\_UPDATE]  \hyperref[regdesc:MEMMONITORLOGMINREG]{MEM\_MONITOR\_LOG\_MIN\_REG} 到 \hyperref[regdesc:MEMMONITORLOGMAXREG]{MEM\_MONITOR\_LOG\_MAX\_REG} 的地址更新为 DMA\_2 总线的监测地址空间。\\
0：不更新\\
1：更新\\
其他值：无效\\
(WT)
\label{fielddesc:MEMMONITORLOGMONADDRDMA3UPDATE}\item [MEM\_MONITOR\_LOG\_MON\_ADDR\_DMA\_3\_UPDATE]  \hyperref[regdesc:MEMMONITORLOGMINREG]{MEM\_MONITOR\_LOG\_MIN\_REG} 到 \hyperref[regdesc:MEMMONITORLOGMAXREG]{MEM\_MONITOR\_LOG\_MAX\_REG} 的地址更新为 DMA\_3 总线的监测地址空间。\\
0：不更新\\
1：更新\\
其他值：无效\\
(WT)
\end{reglist}\end{regdesc}
\end{register}


\begin{register}{H}{MEM\_MONITOR\_LOG\_MEM\_START\_REG}{0x0020}\label{regdesc:MEMMONITORLOGMEMSTARTREG}
\regfield{MEM\_MONITOR\_LOG\_MEM\_START}{32}{0}{{0}}%
\reglabel{Reset}\regnewline%
\begin{regdesc}\begin{reglist}
\label{fielddesc:MEMMONITORLOGMEMSTART}\item [MEM\_MONITOR\_LOG\_MEM\_START] 配置记录数据写入内存的起始地址。(R/W)
\end{reglist}\end{regdesc}
\end{register}


\begin{register}{H}{MEM\_MONITOR\_LOG\_MEM\_END\_REG}{0x0024}\label{regdesc:MEMMONITORLOGMEMENDREG}
\regfield{MEM\_MONITOR\_LOG\_MEM\_END}{32}{0}{{0}}%
\reglabel{Reset}\regnewline%
\begin{regdesc}\begin{reglist}
\label{fielddesc:MEMMONITORLOGMEMEND}\item [MEM\_MONITOR\_LOG\_MEM\_END] 配置记录数据写入内存的结束地址。(R/W)
\end{reglist}\end{regdesc}
\end{register}


\begin{register}{H}{MEM\_MONITOR\_LOG\_MEM\_CURRENT\_ADDR\_REG}{0x0028}\label{regdesc:MEMMONITORLOGMEMCURRENTADDRREG}
\regfield{MEM\_MONITOR\_LOG\_MEM\_CURRENT\_ADDR}{32}{0}{{0}}%
\reglabel{Reset}\regnewline%
\begin{regdesc}\begin{reglist}
\label{fielddesc:MEMMONITORLOGMEMCURRENTADDR}\item [MEM\_MONITOR\_LOG\_MEM\_CURRENT\_ADDR] 表示下一次写内存的写地址。(RO)
\end{reglist}\end{regdesc}
\end{register}


\begin{register}{H}{MEM\_MONITOR\_LOG\_MEM\_ADDR\_UPDATE\_REG}{0x002C}\label{regdesc:MEMMONITORLOGMEMADDRUPDATEREG}
\regfield{\rotatebox{10}{(reserved)}}{31}{1}{0000000000000000000000000000000}%
\regfield{\rotatebox{10}{MEM\_MONITOR\_LOG\_MEM\_ADDR\_UPDATE}}{1}{0}{{0}}%
\reglabel{Reset}\regnewline%
\begin{regdesc}\begin{reglist}
\label{fielddesc:MEMMONITORLOGMEMADDRUPDATE}\item [MEM\_MONITOR\_LOG\_MEM\_ADDR\_UPDATE] 配置是否将 \hyperref[regdesc:MEMMONITORLOGMEMSTARTREG]{MEM\_MONITOR\_LOG\_MEM\_START\_REG} 中的值更新为 \hyperref[regdesc:MEMMONITORLOGMEMCURRENTADDRREG]{MEM\_MONITOR\_LOG\_MEM\_CURRENT\_ADDR\_REG} 中的值。\\
0：不更新\\
1：更新\\
(WT)
\end{reglist}\end{regdesc}
\end{register}


\begin{register}{H}{MEM\_MONITOR\_LOG\_MEM\_FULL\_FLAG\_REG}{0x0030}\label{regdesc:MEMMONITORLOGMEMFULLFLAGREG}
\regfield{\rotatebox{10}{(reserved)}}{30}{2}{000000000000000000000000000000}%
\regfield{\rotatebox{10}{MEM\_MONITOR\_CLR\_LOG\_MEM\_FULL\_FLAG}}{1}{1}{{0}}%
\regfield{\rotatebox{10}{MEM\_MONITOR\_LOG\_MEM\_FULL\_FLAG}}{1}{0}{{0}}%
\reglabel{Reset}\regnewline%
\begin{regdesc}\begin{reglist}
\label{fielddesc:MEMMONITORLOGMEMFULLFLAG}\item [MEM\_MONITOR\_LOG\_MEM\_FULL\_FLAG] 表示数据是否溢出存储地址范围。\\
0：未溢出\\
1：溢出\\
(RO)
\label{fielddesc:MEMMONITORCLRLOGMEMFULLFLAG}\item [MEM\_MONITOR\_CLR\_LOG\_MEM\_FULL\_FLAG] 配置是否清除 \hyperref[fielddesc:MEMMONITORLOGMEMFULLFLAG]{MEM\_MONITOR\_LOG\_MEM\_FULL\_FLAG} 标志位。\\
0：不清除（默认）\\
1：清除\\
(WT)
\end{reglist}\end{regdesc}
\end{register}


\begin{register}{H}{MEM\_MONITOR\_CLOCK\_GATE\_REG}{0x0034}\label{regdesc:MEMMONITORCLOCKGATEREG}
\regfield{(reserved)}{31}{1}{0000000000000000000000000000000}%
\regfield{MEM\_MONITOR\_CLK\_EN}{1}{0}{{0}}%
\reglabel{Reset}\regnewline%
\begin{regdesc}\begin{reglist}
\label{fielddesc:MEMMONITORCLKEN}\item [MEM\_MONITOR\_CLK\_EN] 配置是否使能寄存器时钟门控。\\
0：不使能\\
1：使能\\
(R/W)
\end{reglist}\end{regdesc}
\end{register}


\begin{register}{H}{MEM\_MONITOR\_DATE\_REG}{0x03FC}\label{regdesc:MEMMONITORDATEREG}
\regfield{(reserved)}{4}{28}{0000}%
\regfield{MEM\_MONITOR\_DATE}{28}{0}{{0x2308140}}%
\reglabel{Reset}\regnewline%
\begin{regdesc}\begin{reglist}
\label{fielddesc:MEMMONITORDATE}\item [MEM\_MONITOR\_DATE] 版本控制寄存器。(R/W)
\end{reglist}\end{regdesc}
\end{register}
